\subsection{Measurement under incomplete observability}
\label{subsec:measurement_contribution}

The paper's measurement contribution is to make industrial AI diffusion observable under incomplete, uneven, and administratively mediated data. Industrial AI adoption is difficult to measure because many deployments occur inside production lines, enterprise software, quality-control systems, robotics cells, warehouse systems, and process-control routines. These deployments may affect production without appearing in standard innovation data, firm surveys, or patent records. The empirical strategy therefore uses administrative recognition as an adoption trace, while explicitly separating what is observed, what is inferred, what is externally verified, and what remains outside the data.

The first measurement problem is the distinction between adoption and recognized adoption. Smart-factory lists identify projects that have entered a formal category of intelligent manufacturing, which makes them valuable for studying policy-legible industrial upgrading. These lists do not measure every use of AI inside factories. The paper therefore treats recognition as a trace of visible adoption rather than as a complete census. This distinction protects the empirical interpretation and turns the smart-factory lists into a disciplined measurement object instead of an overextended proxy.

The second measurement problem is spatial assignment. The paper studies diffusion architecture at the city level, so project geography must be resolved with evidence that can be inspected and audited. The city-resolution pipeline assigns every project to a city and classifies the supporting evidence into official-location exact assignments, rule-based text inferences, and externally verified assignments. This tiering matters because different rows carry different degrees of evidentiary strength. The paper does not collapse them into a single unqualified location variable. It reports the evidence classes and uses a stratified audit to make location uncertainty visible.

The third measurement problem is claim discipline. Empirical AI research often combines administrative data, scraped records, inferred geography, patent records, and policy labels, then presents them as if they supported a single level of inference. This paper instead organizes the evidence into claim tiers. Dataset counts and evidence splits are measured facts. Pilot-zone and smart-factory overlap is a descriptive association. Hub attenuation is robustness evidence for a spatial architecture. Sectoral exposure and export relevance are mechanism and context evidence. Public controls and tiered patent geography are appendix robustness exercises. Causal treatment effects, productivity effects, exact patent geocoding, and strict production-control estimates remain outside the completed evidence base.

The fourth measurement problem is reproducibility. The paper is accompanied by a pipeline that regenerates the main submission package, validates the evidence gates, exports paper tables and figures, and blocks unsupported claims. This matters because the empirical object is constructed from multiple administrative and derived sources. Reproducibility is therefore not only a software standard. It is part of the measurement design, since it allows the reader to distinguish stable paper-facing artifacts from exploratory layers, blocked specifications, and robustness-only extensions.

The resulting approach can be used beyond the Chinese case. Many AI-adoption settings will have incomplete observability, because deployment often appears inside workflows and production systems before it appears in public statistics. A general measurement strategy for AI diffusion should therefore combine visible administrative traces, evidence-tiered geography, audit procedures, claim-specific robustness, and transparent limits on causal interpretation. The paper contributes such a strategy by showing how an adoption architecture can be measured without pretending that the observable trace is the full adoption universe.